% Figure générée par Figurine -- format figurine/1, thème « these ».
% Nécessite \usepackage{figurine} (figurine.sty, copiable depuis l'appli).
% Unités : mm ; origine en bas à gauche de la planche (80 × 40 mm).
\begin{tikzpicture}[figurine]
\useasboundingbox (0,0) rectangle (80,40);

% a : Rectangle -- 20 × 12 mm, « Bloc »
\draw[fig/trait] (5,23) rectangle (25,35);
\node[fig/texte, anchor=center] at (15,29) {Bloc};

% b : Rectangle -- 20 × 12 mm, « $E_1$ »
\filldraw[fig/trait, fig/fond gris] (30,23) rectangle (50,35);
\node[fig/texte, anchor=center] at (40,29) {$E_1$};

% c : Rectangle -- 20 × 12 mm, « 50 % & plus »
\filldraw[fig/trait, fig/hachure bitumineux] (55,23) rectangle (75,35);
\node[fig/etiquette, anchor=center] at (65,29) {50 \% \& plus};

% d : Rectangle -- 16 × 6 mm
\filldraw[fig/trait, fig/hachure granulaire] (20,6) -- (33.86,14) -- (36.86,8.8) -- (23,0.8) -- cycle;

% e : Rectangle -- 10 × 5 mm
\filldraw[fig/trait, fig/hachure sol] (55,18) rectangle (65,23);
\end{tikzpicture}
